\chapter{Results}
\label{ch:results}

This chapter presents the experimental results from Stages~3 and~4. Stage~3
establishes a text-only generative baseline (Mistral-7B); Stage~4 introduces
layout-aware discriminative extraction (LayoutLMv3 + Biaffine Linker) in two
linking variants: token-level (V1) and span-level (V2).

% ===========================================================================
\section{Dataset Coverage}
\label{sec:dataset-coverage}

The evaluation was performed on the KVP10k test set. The original HuggingFace
test split contains 5{,}255 annotated rows. After annotator deduplication
(approximately 5 copies per unique page) and filtering for pages successfully
prepared with PyMuPDF, the final evaluation set contains \textbf{581 unique
test pages} ($\sim$55\% of the 1{,}051 unique pages). This coverage is
limited by PDF download failures and scanned PDFs lacking usable text layers.

% ===========================================================================
\section{Stage~3: Mistral-7B Generative Baseline}
\label{sec:stage3-results}

\subsection{Overall Performance}

Table~\ref{tab:stage3-overall} summarises the Mistral-7B baseline under two
threshold configurations.

\begin{table}[ht]
  \centering
  \caption{Stage~3 (Mistral-7B QLoRA) overall results ($n = 581$ test pages).}
  \label{tab:stage3-overall}
  \begin{tabular}{@{}lcccc@{}}
    \toprule
    \textbf{Protocol} & \textbf{Precision} & \textbf{Recall} & \textbf{F1} & \textbf{TP / FP / FN} \\
    \midrule
    Text-only (NED$<$0.2) & 0.604 & 0.755 & 0.671 & 9{,}482 / 6{,}225 / 3{,}070 \\
    Text+Bbox (NED$<$0.2, IoU$>$0.3) & 0.539 & 0.675 & 0.600 & 8{,}473 / 7{,}234 / 4{,}079 \\
    \bottomrule
  \end{tabular}
\end{table}

\subsection{Paper Baseline Comparison}

Table~\ref{tab:paper-comparison} compares our Mistral baseline against the
published KVP10k baseline~\cite{kvp10k}.

\begin{table}[ht]
  \centering
  \caption{Published KVP10k baseline vs.\ our Stage~3 Mistral baseline.}
  \label{tab:paper-comparison}
  \begin{tabular}{@{}lcc@{}}
    \toprule
    \textbf{System} & \textbf{Text-only F1} & \textbf{Text+Location F1} \\
    \midrule
    KVP10k paper baseline (Table~1, All) & 0.643 & 0.612 \\
    Ours (NED$<$0.2, IoU$>$0.3; $n=581$) & 0.671 & 0.600 \\
    \bottomrule
  \end{tabular}
\end{table}

Our text-only F1 is higher than the published baseline, but text+location F1
is slightly lower. This comparison is approximate because our evaluation
covers only 55\% of unique test pages due to PDF availability constraints.

\subsection{Per-Cluster Analysis}
\label{sec:stage3-per-cluster}

Table~\ref{tab:stage3-cluster} breaks down Stage~3 performance by layout
cluster, revealing that Mistral's errors are heavily layout-dependent.

\begin{table}[ht]
  \centering
  \caption{Stage~3 Mistral error analysis by layout cluster ($n = 581$).}
  \label{tab:stage3-cluster}
  \begin{tabular}{@{}lcccccc@{}}
    \toprule
    \textbf{Cluster} & \textbf{Docs} & \textbf{Pred} & \textbf{GT}
                     & \textbf{Correct (\%)} & \textbf{Halluc.\ (\%)}
                     & \textbf{Missed (\%)} \\
    \midrule
    0 (Dense)  & 465 & 7{,}852 & 8{,}161 & 48.7 & 40.5 & 31.2 \\
    1 (Sparse) & 116 &    889 &    224 & 14.4 & 52.9 & 41.1 \\
    \bottomrule
  \end{tabular}
\end{table}

Mistral achieves a $3.4\times$ higher correct-match rate on dense
layouts. Sparse layouts suffer from catastrophic over-prediction (889
predictions vs.\ 224 GT, a 3.97$\times$ over-generation ratio) and a
hallucination rate of 52.9\%. This asymmetry motivates Stage~4: layout-aware
models should reduce hallucination by grounding predictions in spatial context.

% ===========================================================================
\section{Stage~4: LayoutLMv3 Discriminative Models}
\label{sec:stage4-results}

\subsection{Stage~4a: Entity Classification}

Stage~4a trains the LayoutLMv3 encoder with the entity classification head
only ($\lambda = 0$). The best model achieves:

\begin{table}[ht]
  \centering
  \caption{Stage~4a entity classification results (best epoch).}
  \label{tab:stage4a-results}
  \begin{tabular}{@{}lccc@{}}
    \toprule
    \textbf{Metric} & \textbf{Precision} & \textbf{Recall} & \textbf{F1} \\
    \midrule
    Entity (Key + Value) & --- & --- & 0.847 \\
    \bottomrule
  \end{tabular}
\end{table}

Entity F1 of 0.847 confirms that the discriminative LayoutLMv3 model
effectively classifies tokens as keys, values, or other --- establishing
strong entity representations for the linking stage.

\subsection{Stage~4b V1: Token-Level Linking}
\label{sec:v1-results}

Stage~4b V1 activates the biaffine linker at the token level. After full
training (12 epochs, best model at epoch~9), the results are:

\begin{table}[ht]
  \centering
  \caption{Stage~4b V1 (token-level linking) results.}
  \label{tab:v1-results}
  \begin{tabular}{@{}lcccc@{}}
    \toprule
    \textbf{Metric} & \textbf{Precision} & \textbf{Recall} & \textbf{F1} & \textbf{TP / Pred / GT} \\
    \midrule
    Entity F1 & --- & --- & 0.841 & --- \\
    Link F1   & 0.075 & 0.004 & 0.008 & 15 / 200 / 3{,}776 \\
    \bottomrule
  \end{tabular}
\end{table}

\paragraph{Interpretation.}
Token-level linking achieves near-zero link F1 despite reasonable entity
classification. The model predicts only 200 links (of which 15 are correct)
against 3{,}776 ground-truth links. The root cause is the extreme class
imbalance at the token level: on a typical page, $\sim$100 key tokens
$\times$ $\sim$100 value tokens produces $\sim$10{,}000 candidate pairs with
only $\sim$20 positives ($\sim$450:1 neg:pos ratio). Even with
\texttt{pos\_weight} = 50, the linker cannot overcome this imbalance and
becomes overly conservative, predicting very few links.

\paragraph{Lambda sweep.}
The $\lambda$ sweep over $\{0.5, 1.0, 2.0\}$ shows that entity F1 is robust
to the linking weight ($\lambda = 0.5$: F1 = 0.849; $\lambda = 1.0$: F1 = 0.845;
$\lambda = 2.0$: F1 = 0.848) but link F1 remains near zero across all
values. This confirms that the failure is structural (granularity mismatch),
not a hyperparameter issue.

\subsection{Stage~4b V2: Span-Level Linking}
\label{sec:v2-results}

Stage~4b V2 addresses the token-level granularity mismatch by grouping
contiguous same-label tokens into spans before linking
(Section~\ref{sec:v2-span-linking}). This reduces the candidate matrix from
$\sim$10{,}000 to $\sim$225 entries with a $\sim$5:1 neg:pos ratio.

The best V2 model (trained with teacher forcing, early stopped at epoch~12
with patience~7) achieves:

\begin{table}[ht]
  \centering
  \caption{Stage~4b V2 (span-level linking) results.}
  \label{tab:v2-results}
  \begin{tabular}{@{}lcccc@{}}
    \toprule
    \textbf{Metric} & \textbf{Precision} & \textbf{Recall} & \textbf{F1} & \textbf{TP / Pred / GT} \\
    \midrule
    Entity F1 & --- & --- & 0.847 & --- \\
    Link F1   & 0.302 & 0.003 & 0.007 & 37 / 498 / 3{,}776 \\
    \bottomrule
  \end{tabular}
\end{table}

\paragraph{Interpretation.}
Despite the favourable neg:pos ratio at the span level, V2 link F1 remains
near zero.  Three root causes were identified through the diagnostic
experiments described in Section~\ref{sec:diagnostic-results}.

% ===========================================================================
\section{Stage~3 vs.\ Stage~4 Comparison}
\label{sec:comparison}

Table~\ref{tab:overall-comparison} summarises the progression from generative
baseline to discriminative models.

\begin{table}[ht]
  \centering
  \caption{Overall comparison across all stages. Dashes indicate metrics not
    applicable to that system: the KVP10k baseline and Mistral produce
    end-to-end KVP predictions without a separate entity-classification head.
    Mistral's F1 is text\,+\,location (NED\,$<$\,0.2, IoU\,$>$\,0.3).}
  \label{tab:overall-comparison}
  \small
  \begin{tabular}{@{}lccp{5.5cm}@{}}
    \toprule
    \textbf{Model} & \textbf{Entity F1} & \textbf{Link F1} & \textbf{Key Finding} \\
    \midrule
    KVP10k paper baseline     & ---   & 0.659 & Official text-only generative system \\
    Stage~3 Mistral (text+loc) & ---  & 0.600 & 52.9\% hallucination on sparse layouts \\
    Stage~4a (entity only)     & 0.847 & ---   & Strong entity classification established \\
    Stage~4b V1 (token linking) & 0.841 & 0.008 & Token-level class imbalance ($\sim$450:1) kills linking \\
    Stage~4b V2 (span linking)  & 0.847 & 0.007 & Data pipeline bugs mask span-level benefit \\
    Stage~4b V4 (corrected)     & 0.827 & 0.673 & Pipeline fixes recover linking; matches KVP10k baseline \\
    \bottomrule
  \end{tabular}
\end{table}

% ===========================================================================
\section{Ablation: Linking Granularity}
\label{sec:linking-ablation}

The comparison between V1 and V2 constitutes the central ablation of this
thesis. The key metrics are:

\begin{table}[ht]
  \centering
  \caption{V1 (token-level) vs.\ V2 (span-level) vs.\ V4 (corrected) linking ablation.}
  \label{tab:v1-v2-ablation}
  \small
  \begin{tabular}{@{}lcccc@{}}
    \toprule
    \textbf{Variant} & \textbf{Cand./page} & \textbf{Neg:Pos}
                     & \textbf{Link F1} & \textbf{Entity F1} \\
    \midrule
    V1 (token) & $\sim$10{,}000 & $\sim$450:1 & 0.008 & 0.841 \\
    V2 (span)  & $\sim$225 & $\sim$5:1 & 0.007 & 0.847 \\
    V4 (corrected) & $\sim$225 & $\sim$5:1 & 0.673 & 0.827 \\
    \bottomrule
  \end{tabular}
\end{table}

% ===========================================================================
\section{Diagnostic Findings}
\label{sec:diagnostic-results}

The near-zero link F1 of V2 prompted a systematic investigation of the
training pipeline.  Three root causes were identified:

\subsection{Bug~1: Missing Teacher Forcing}

During V2 training, the linker received \texttt{argmax} predictions from the
entity classifier rather than ground-truth labels for span grouping.  This
means the linker's span boundaries were corrupted by entity classification
errors, and the link labels (aligned to ground-truth spans) were misaligned
with the linker's input.  The fix applies teacher forcing during training,
where the linker receives ground-truth entity labels.

After applying teacher forcing, the best V2+TF model achieved entity
F1\,=\,0.847 and link F1\,=\,0.007 --- unchanged, indicating deeper issues.

\subsection{Bug~2: Bounding Box Filter}

The \texttt{group\_contiguous\_spans} function used the predicate
$(x_1 < x_2) \wedge (y_1 < y_2)$ to filter out invalid bounding boxes.
However, LayoutLMv3's normalisation to the $[0, 1000]$ grid often produces
coordinates where $x_1 = x_2$ (zero-width tokens).  The filter rejected
\textbf{49\%} of predicted key tokens and \textbf{46\%} of predicted value
tokens, cutting the number of entity spans roughly in half.  Only
$[0,0,0,0]$ bounding boxes (CLS, SEP, PAD tokens) should be excluded.

\subsection{Bug~3: Link Label Cross-Product Explosion}

The label generation function \texttt{\_generate\_labels} created a link
entry for \emph{every} key-word index $\times$ \emph{every} value-word
index within a KVP.  For a typical document with 8~KVPs, this produced
$\sim$1{,}049 positive entries in the $512 \times 512$ link label matrix
instead of~8.  After span-level collapse, nearly every key--value span pair
was marked positive, providing no discrimination for the linker to learn
from.

Compounding this, the OCR text blocks were multi-word phrases stored as
single ``words'' (Section~\ref{sec:text-processing}), so the
bounding-box overlap matching frequently mapped multiple KVPs to the same
block index, further inflating the cross-product.

\subsection{V4 Corrections and Verification}

Table~\ref{tab:v4-verification} summarises the effect of the four V4
corrections on 20 test documents.

\begin{table}[ht]
  \centering
  \caption{Link label pipeline verification (20 test documents, 179 GT KVPs).}
  \label{tab:v4-verification}
  \begin{tabular}{@{}lcc@{}}
    \toprule
    \textbf{Metric} & \textbf{Before (V2)} & \textbf{After (V4)} \\
    \midrule
    Words per document & $\sim$35 & $\sim$200 \\
    Token-level link positives (20 docs) & $\sim$20{,}000 & 1{,}013 \\
    Word-level link recovery & 100\%\textsuperscript{*} & 79\% (141/179) \\
    Span-level link recovery & 100\%\textsuperscript{*} & 70\% (125/179) \\
    \bottomrule
  \end{tabular}
\end{table}

\noindent\textsuperscript{*}\,Before V4, ``100\% recovery'' is misleading:
the cross-product produced \emph{all} possible positive links, not just
correct ones.  The effective supervision was a near-uniform positive matrix,
equivalent to no supervision.

% ===========================================================================
\section{Computational Performance}
\label{sec:computational}

\begin{table}[ht]
  \centering
  \caption{Training time and resource usage.}
  \label{tab:compute}
  \begin{tabular}{@{}lccc@{}}
    \toprule
    \textbf{Stage} & \textbf{GPU} & \textbf{Time/epoch} & \textbf{Total epochs} \\
    \midrule
    Stage~3 (Mistral QLoRA) & A100 40\,GB & $\sim$45 min & 8 \\
    Stage~4a (entity) & A100 40\,GB & $\sim$1.5 hrs & 20 \\
    Stage~4b V1 (token link) & A100 40\,GB & $\sim$1.8 hrs & 12 (best: 9) \\
    Stage~4b V2 (span link) & A100 40\,GB & $\sim$1.8 hrs & 12 (best: epoch~12) \\
    Stage~4b V4 (corrected) & A100 40\,GB & $\sim$1.8 hrs & 30 (best: 13) \\
    \bottomrule
  \end{tabular}
\end{table}
